\section*{Introduction}
Nervous system disorders are the leading cause of disability worldwide and account for a growing share of healthcare costs in high-income countries. As of 2025, they affected more than 40\% of the global population \href{https://www.who.int/publications/i/item/9789240116139}{[1]}. Yet new treatments remain scarce. Only 5.9\% of neurological and 7.3\% of psychiatric drug programs entering Phase 1 reach approval, below the 7.9\% average across all indications \href{https://www.bio.org/clinical-development-success-rates-and-contributing-factors-2011-2020}{[2]}. Many of these setbacks reflect gaps in our understanding of neurobiology. On the forefront are the researchers that face the challenge of large amount of data from cell and animal models that poorly recapitulate human disease, drug mechanisms that cannot produce a sufficient therapeutic effect, and clinical endpoints misaligned with disease heterogeneity and patient diversity.

Data acquisition is arguably no longer the main constraint on closing these gaps. Brain-wide cell atlases, regional and whole-brain connectomes, biobanks, genomic databases, population imaging studies, and electrophysiology archives have greatly expanded the data available to neuroscientists. The bottleneck has shifted from generating data to analyzing and interpreting it against an immense body of prior knowledge that no single scientist or research team can fully internalize. AI agents working alongside researchers could accelerate this work, but their reliability must be monitored to avoid propagating errors.

Existing benchmarks do not fully test these analytical and reasoning abilities. BrainBench showed that large language models (LLMs) outperform human experts at identifying which of two versions of an abstract reports a study's true results \href{https://doi.org/10.1038/s41562-024-02046-9}{[3]}. Such benchmarks measure how well a model has absorbed the existing literature, not whether it can analyze new data. Agentic biology benchmarks such as LAB-Bench \href{https://arxiv.org/abs/2407.10362}{[4]} and BixBench \href{https://arxiv.org/abs/2503.00096}{[5]}, and our own SpatialBench \href{https://arxiv.org/abs/2512.21907}{[6]}, scBench \href{https://arxiv.org/abs/2602.09063}{[7]}, and EpiBench \href{https://arxiv.org/abs/2606.13602}{[8]}, do test analysis of experimental data, but none is built around the data types specific to neuroscience. Our disease-area benchmarks, TxBench \href{https://arxiv.org/abs/2606.19245}{[9]}\href{https://latch.bio/txbench-ab}{[10]}\href{https://latch.bio/txbench-oligo}{[11]}\href{https://latch.bio/txbench-ab-dev}{[12]}, OncologyBench \href{https://benchmarks.bio/oncology}{[13]}, and RareDiseaseBench \href{https://benchmarks.bio/raredisease}{[14]}, have shown that frontier agents can execute individual analyses competently and the failure is often one of reasoning rather than the actual execution. NeuroscienceBench applies this same design to neuroscience, where diverse data types and scientific decision making requires deep domain expertise.
